\section{反向过程的参数化}

\subsection{反向转移分布的建模}

既然真实后验 $q(x_{t-1}|x_t, x_0)$ 是高斯分布，我们自然将反向转移 $p_\theta(x_{t-1}|x_t)$ 也建模为高斯分布：

$$p_\theta(x_{t-1}|x_t) = \mathcal{N}\big(x_{t-1};\, \mu_\theta(x_t, t),\; \sigma_t^2 I\big)$$

其中：
\begin{itemize}
    \item $\mu_\theta(x_t, t)$：由神经网络预测的均值，以 $x_t$ 和时间步 $t$ 为输入
    \item $\sigma_t^2$：方差，设定为 $\tilde{\beta}_t$（与真实后验相同），\textbf{不需要学习}
\end{itemize}

\begin{knowledgebox}{为什么方差不需要学习？}
由于我们要最小化 $D_{\text{KL}}(q(x_{t-1}|x_t,x_0) \| p_\theta(x_{t-1}|x_t))$，而两个高斯分布之间的 KL 散度在方差相同时简化为均值之差的 L2 范数。因此将 $p_\theta$ 的方差设为与 $q$ 相同的 $\tilde{\beta}_t$，训练目标就只剩下匹配均值。后续的 Improved DDPM 论文探索了学习方差的方案，但基本 DDPM 中固定方差即可。
\end{knowledgebox}

\subsection{简化后的训练目标}

将高斯形式代入 KL 散度，去噪匹配项 $L_{t-1}$ 简化为：

$$L_{t-1} = \frac{1}{2\sigma_t^2} \left\| \tilde{\mu}_t(x_t, x_0) - \mu_\theta(x_t, t) \right\|^2 + C$$

其中 $C$ 是与 $\theta$ 无关的常数。因此，训练目标就是让网络预测的均值 $\mu_\theta(x_t, t)$ 尽可能接近真实后验均值 $\tilde{\mu}_t(x_t, x_0)$。

\begin{importantbox}{训练的直觉}
训练过程可以简洁地描述为：（1）采样一个干净数据 $x_0$；（2）采样时间步 $t$；（3）通过前向跳跃分布采样含噪数据 $x_t$；（4）用神经网络预测均值 $\mu_\theta(x_t, t)$；（5）计算预测均值与真实均值的 L2 损失并反向传播。
\end{importantbox}

实际训练中，系数 $\frac{1}{2\sigma_t^2}$ 相当于一个与时间 $t$ 相关的权重项。Ho et al.（2020）发现，在实践中\textbf{去掉这个权重}（即对所有时间步使用均匀权重）效果也很好，甚至更好。这是 DDPM 论文中的一个重要经验发现。

\subsection{本章小结}

将反向转移建模为高斯分布并固定方差后，训练目标简化为均值的 L2 回归问题。权重项在实践中常被省略以简化训练。

%% ========== 第五节：三种预测目标 ==========

